\documentclass[10pt,a4paper]{amsart}
\usepackage{amsmath,amssymb,amsthm,graphicx}
% For commutative diagrams we use Xy-pic
\usepackage[all]{xy}

% We use 2cell for 2-commutative diagrams.
\xyoption{2cell}
% Theorem environments.
%
\theoremstyle{plain}
\newtheorem{theorem}[subsection]{Theorem}
\newtheorem{proposition}[subsection]{Proposition}
\newtheorem{lemma}[subsection]{Lemma}

\theoremstyle{definition}
\newtheorem{definition}[subsection]{Definition}
\newtheorem{example}[subsection]{Example}
\newtheorem{exercise}[subsection]{Exercise}
\newtheorem{situation}[subsection]{Situation}

\theoremstyle{remark}
\newtheorem{remark}[subsection]{Remark}
\newtheorem{remarks}[subsection]{Remarks}

\numberwithin{equation}{subsection}

% Macros
%
\def\lim{\mathop{\mathrm{lim}}\nolimits}
\def\colim{\mathop{\mathrm{colim}}\nolimits}
\def\Spec{\mathop{\mathrm{Spec}}}
\def\Hom{\mathop{\mathrm{Hom}}\nolimits}
\def\Ext{\mathop{\mathrm{Ext}}\nolimits}
\def\SheafHom{\mathop{\mathcal{H}\!\mathit{om}}\nolimits}
\def\SheafExt{\mathop{\mathcal{E}\!\mathit{xt}}\nolimits}
\def\Sch{\mathit{Sch}}
\def\Mor{\mathop{\mathrm{Mor}}\nolimits}
\def\Ob{\mathop{\mathrm{Ob}}\nolimits}
\def\Sh{\mathop{\mathit{Sh}}\nolimits}
\def\NL{\mathop{N\!L}\nolimits}
\def\CH{\mathop{\mathrm{CH}}\nolimits}
\def\proetale{{pro\text{-}\acute{e}tale}}
\def\etale{{\acute{e}tale}}
\def\QCoh{\mathit{QCoh}}
\def\Ker{\mathop{\mathrm{Ker}}}
\def\Im{\mathop{\mathrm{Im}}}
\def\Coker{\mathop{\mathrm{Coker}}}
\def\Coim{\mathop{\mathrm{Coim}}}

% Boxtimes
%
\DeclareMathSymbol{\boxtimes}{\mathbin}{AMSa}{"02}

%
% Macros for moduli stacks/spaces
%
\def\QCohstack{\mathcal{QC}\!\mathit{oh}}
\def\Cohstack{\mathcal{C}\!\mathit{oh}}
\def\Spacesstack{\mathcal{S}\!\mathit{paces}}
\def\Quotfunctor{\mathrm{Quot}}
\def\Hilbfunctor{\mathrm{Hilb}}
\def\Curvesstack{\mathcal{C}\!\mathit{urves}}
\def\Polarizedstack{\mathcal{P}\!\mathit{olarized}}
\def\Complexesstack{\mathcal{C}\!\mathit{omplexes}}
% \Pic is the operator that assigns to X its picard group, usage \Pic(X)
% \Picardstack_{X/B} denotes the Picard stack of X over B
% \Picardfunctor_{X/B} denotes the Picard functor of X over B
\def\Pic{\mathop{\mathrm{Pic}}\nolimits}
\def\Picardstack{\mathcal{P}\!\mathit{ic}}
\def\Picardfunctor{\mathrm{Pic}}
\def\Deformationcategory{\mathcal{D}\!\mathit{ef}}
\usepackage{cleveref}
\crefname{originallemma}{original lemma}{original lemmas}
\Crefname{originallemma}{Original lemma}{Original lemmas}
\makeatletter
\newlabel{lemma-pullback-constructible}{{095G}{source}{}{sourcepointer-1}{}}
\newlabel{lemma-pullback-constructible@cref}{{[originallemma][0][]095G}{[1][0][]source}{}{}{}}
\makeatother
\crefname{originallemma}{original lemma}{original lemmas}
\Crefname{originallemma}{Original lemma}{Original lemmas}
\makeatletter
\newlabel{lemma-torsion-colimit-constructible}{{03SA}{source}{}{sourcepointer-2}{}}
\newlabel{lemma-torsion-colimit-constructible@cref}{{[originallemma][0][]03SA}{[1][0][]source}{}{}{}}
\makeatother
\crefname{originallemma}{original lemma}{original lemmas}
\Crefname{originallemma}{Original lemma}{Original lemmas}
\makeatletter
\newlabel{categories-lemma-adjoint-exact}{{0038}{source}{}{sourcepointer-3}{}}
\newlabel{categories-lemma-adjoint-exact@cref}{{[originallemma][0][]0038}{[1][0][]source}{}{}{}}
\makeatother
\crefname{originallemma}{original lemma}{original lemmas}
\Crefname{originallemma}{Original lemma}{Original lemmas}
\makeatletter
\newlabel{lemma-colimit-constructible}{{095P}{source}{}{sourcepointer-4}{}}
\newlabel{lemma-colimit-constructible@cref}{{[originallemma][0][]095P}{[1][0][]source}{}{}{}}
\makeatother
\crefname{originallemma}{original lemma}{original lemmas}
\Crefname{originallemma}{Original lemma}{Original lemmas}
\makeatletter
\newlabel{lemma-stalk-pullback}{{03Q1}{source}{}{sourcepointer-5}{}}
\newlabel{lemma-stalk-pullback@cref}{{[originallemma][0][]03Q1}{[1][0][]source}{}{}{}}
\makeatother
\crefname{originaltheorem}{original theorem}{original theorems}
\Crefname{originaltheorem}{Original theorem}{Original theorems}
\makeatletter
\newlabel{theorem-exactness-stalks}{{03PU}{source}{}{sourcepointer-6}{}}
\newlabel{theorem-exactness-stalks@cref}{{[originaltheorem][0][]03PU}{[1][0][]source}{}{}{}}
\makeatother
\crefname{originallemma}{original lemma}{original lemmas}
\Crefname{originallemma}{Original lemma}{Original lemmas}
\makeatletter
\newlabel{lemma-constructible-abelian}{{03RZ}{source}{}{sourcepointer-7}{}}
\newlabel{lemma-constructible-abelian@cref}{{[originallemma][0][]03RZ}{[1][0][]source}{}{}{}}
\makeatother
\title{Constructibility can be tested after a surjective morphism of quasi-compact quasi-separated schemes}
\author{The Stacks Project Authors}
\date{stacks-project:etale-cohomology.tex:local-snapshot-20261011; GNU FDL 1.2 or later}
\begin{document}
\maketitle
\noindent\textit{Editorial note.} All delivered mathematical text is the author's own, copied from the pinned Stacks Project chapter file etale-cohomology.tex; no statement or proof step was rewritten, reordered or removed.\par\medskip
\noindent\textit{Editorial note.} The proof is the author's complete argument for all three parts; the two directions are both written out.\par\medskip
\noindent\textit{Editorial note.} Every labelled reference occurring inside the counted spans is carried in assembly.external\_references as a non-load-bearing navigation pointer (same-chapter ones at their definition site, cross-chapter ones at the invocation inside the counted span), so no load-bearing mathematics is replaced by a stub. No bibliography citation occurs inside the counted spans.\par\medskip
\noindent\textit{Editorial note.} The checked chapter list for same-name collisions is the set of stacks chapters staged in this repository: chow, cohomology, divisors, etale-cohomology, flat, more-algebra, more-morphisms, morphisms, spaces-more-morphisms, varieties together with this card's chapter etale-cohomology.\par\medskip
\noindent\textit{Editorial note.} The preamble is the chapter's own native preamble excerpt already used for this source: the Xy-pic setup (native preamble lines 14-18) and the native theorem declarations with the mathematical macros (native preamble lines 154-217).\par\medskip
\noindent\textit{Allowed standard prerequisites (editorial).}\begin{itemize}
\item Constructible sheaf of sets, of abelian groups or of modules on the small etale site. Same chapter, the definition of constructibility by a partition into locally closed constructible strata on which the sheaf is locally constant. Supplies the property tested in the three parts of the statement.
\item Pullback of a constructible sheaf along a morphism of finite type and the converse for a surjective morphism. Same chapter, the stability of constructibility under pullback, and the criterion producing a partition of Y from one of X along a surjective morphism of quasi-compact quasi-separated schemes. Supplies both directions of the equivalence proved in the proof.
\end{itemize}\medskip
\ifdefined\hypertarget\hypertarget{sourcepointer-1}{}\fi
\noindent\textit{Source reference.} Original lemma 095G: sources/stacks-etale-cohomology/pinned-original-source.tex, source lines 11808--11808. Navigation pointer invoked at proof line 12292; the counted span contains the pointer only and the cited result is not reproduced inside it, so the pointer is registered as a non-load-bearing navigation stub in the established form for Stacks material. No bibliography citation is involved.\par\medskip
\ifdefined\hypertarget\hypertarget{sourcepointer-2}{}\fi
\noindent\textit{Source reference.} Original lemma 03SA: sources/stacks-etale-cohomology/pinned-original-source.tex, source lines 12216--12216. Navigation pointer invoked at proof line 12296; the counted span contains the pointer only and the cited result is not reproduced inside it, so the pointer is registered as a non-load-bearing navigation stub in the established form for Stacks material. No bibliography citation is involved.\par\medskip
\ifdefined\hypertarget\hypertarget{sourcepointer-3}{}\fi
\noindent\textit{Source reference.} Original lemma 0038: sources/stacks-etale-cohomology/pinned-original-source.tex, source lines 12298--12298. Navigation pointer invoked at proof line 12298; the counted span contains the pointer only and the cited result is not reproduced inside it, so the pointer is registered as a non-load-bearing navigation stub in the established form for Stacks material. No bibliography citation is involved.\par\medskip
\ifdefined\hypertarget\hypertarget{sourcepointer-4}{}\fi
\noindent\textit{Source reference.} Original lemma 095P: sources/stacks-etale-cohomology/pinned-original-source.tex, source lines 11903--11903. Navigation pointer invoked at proof line 12300; the counted span contains the pointer only and the cited result is not reproduced inside it, so the pointer is registered as a non-load-bearing navigation stub in the established form for Stacks material. No bibliography citation is involved.\par\medskip
\ifdefined\hypertarget\hypertarget{sourcepointer-5}{}\fi
\noindent\textit{Source reference.} Original lemma 03Q1: sources/stacks-etale-cohomology/pinned-original-source.tex, source lines 4649--4649. Navigation pointer invoked at proof line 12304; the counted span contains the pointer only and the cited result is not reproduced inside it, so the pointer is registered as a non-load-bearing navigation stub in the established form for Stacks material. No bibliography citation is involved.\par\medskip
\ifdefined\hypertarget\hypertarget{sourcepointer-6}{}\fi
\noindent\textit{Source reference.} Original theorem 03PU: sources/stacks-etale-cohomology/pinned-original-source.tex, source lines 3546--3546. Navigation pointer invoked at proof line 12305; the counted span contains the pointer only and the cited result is not reproduced inside it, so the pointer is registered as a non-load-bearing navigation stub in the established form for Stacks material. No bibliography citation is involved.\par\medskip
\ifdefined\hypertarget\hypertarget{sourcepointer-7}{}\fi
\noindent\textit{Source reference.} Original lemma 03RZ: sources/stacks-etale-cohomology/pinned-original-source.tex, source lines 11828--11828. Navigation pointer invoked at proof line 12247; the counted span contains the pointer only and the cited result is not reproduced inside it, so the pointer is registered as a non-load-bearing navigation stub in the established form for Stacks material. No bibliography citation is involved.\par\medskip
\begin{lemma}
\label{lemma-check-constructible}
Let $f : X \to Y$ be a surjective morphism of quasi-compact and
quasi-separated schemes.
\begin{enumerate}
\item Let $\mathcal{F}$ be a sheaf of sets on $Y_\etale$. Then $\mathcal{F}$
is constructible if and only if $f^{-1}\mathcal{F}$ is constructible.
\item Let $\mathcal{F}$ be an abelian sheaf on $Y_\etale$. Then $\mathcal{F}$
is constructible if and only if $f^{-1}\mathcal{F}$ is constructible.
\item Let $\Lambda$ be a Noetherian ring.
Let $\mathcal{F}$ be sheaf of $\Lambda$-modules on $Y_\etale$.
Then $\mathcal{F}$ is constructible if and only if $f^{-1}\mathcal{F}$
is constructible.
\end{enumerate}
\end{lemma}

\begin{proof}
One implication follows from Lemma \ref{lemma-pullback-constructible}.
For the converse, assume $f^{-1}\mathcal{F}$ is constructible.
Write $\mathcal{F} = \colim \mathcal{F}_i$ as a
filtered colimit of constructible sheaves (of sets, abelian groups, or modules)
using Lemma \ref{lemma-torsion-colimit-constructible}.
Since $f^{-1}$ is a left adjoint it commutes with colimits
(Categories, Lemma \ref{categories-lemma-adjoint-exact}) and we see that
$f^{-1}\mathcal{F} = \colim f^{-1}\mathcal{F}_i$.
By Lemma \ref{lemma-colimit-constructible} we see that
$f^{-1}\mathcal{F}_i \to f^{-1}\mathcal{F}$
is surjective for all $i$ large enough.
Since $f$ is surjective we conclude (by looking at stalks using
Lemma \ref{lemma-stalk-pullback} and
Theorem \ref{theorem-exactness-stalks})
that $\mathcal{F}_i \to \mathcal{F}$ is surjective for all $i$ large enough.
Thus $\mathcal{F}$ is the quotient of a constructible sheaf $\mathcal{G}$.
Applying the argument once more to
$\mathcal{G} \times_\mathcal{F} \mathcal{G}$ or
the kernel of $\mathcal{G} \to \mathcal{F}$
we conclude using that $f^{-1}$ is exact and that the category of
constructible sheaves (of sets, abelian groups, or modules) is
preserved under finite (co)limits or (co)kernels inside
$\Sh(Y_\etale)$, $\Sh(X_\etale)$, $\textit{Ab}(Y_\etale)$,
$\textit{Ab}(X_\etale)$, $\textit{Mod}(Y_\etale, \Lambda)$, and
$\textit{Mod}(X_\etale, \Lambda)$, see
Lemma \ref{lemma-constructible-abelian}.
\end{proof}

\end{document}
